\chapter{Tile Framebuffer, Double Buffering, and Hardware Z-Buffer}
\label{chap:tile_memory}

\section{Memory Architecture Overview}
Each of the ten horizontal tiles is backed by an independent on-chip memory subsystem implemented in \texttt{tile.vhd}. The module integrates:
\begin{itemize}
    \item A \textbf{Color VRAM} organized as a True Dual-Port Block RAM supporting simultaneous write operations at $100.0\,\text{MHz}$ and display scan read operations at $25.0\,\text{MHz}$.
    \item A \textbf{Ping-Pong Double Buffering} mechanism that isolates the active rendering target from the scanning display.
    \item An integrated \textbf{Hardware Z-Buffer (Depth Buffer)} providing single-cycle depth testing per pixel.
    \item A \textbf{Hardware Auto-Clear Engine} that resets the back buffer to background white and resets depth to maximum distance with zero CPU cycle consumption.
    \item A \textbf{Multiplier-Free Address Generator} utilizing arithmetic bit shifts in place of DSP multipliers.
\end{itemize}

Figure~\ref{fig:tile_memory} illustrates the complete data-path and control architecture of \texttt{tile.vhd}.

\begin{figure}[htbp]
\centering
\resizebox{0.95\textwidth}{!}{
    \begin{tikzpicture}[
    >=Stealth,
    node distance=1.0cm and 1.2cm,
    every text node part/.style={align=center},
    box/.style={draw=black!80, thick, rounded corners=3pt, fill=white, inner sep=5pt, font=\small},
    membox/.style={box, fill=purple!10, draw=purple!80, minimum width=3.4cm, minimum height=1.6cm},
    logicbox/.style={box, fill=cyan!10, draw=cyan!80, minimum width=2.8cm},
    fsmbox/.style={box, fill=amber!15, draw=amber!80, minimum width=3.2cm},
    bus/.style={draw=black!70, line width=1.3pt, ->},
    sig/.style={draw=blue!70!black, thick, ->},
    ctl/.style={draw=red!80!black, dashed, thick, ->}
]

% ==================== INPUT WRITE SIGNALS ====================
\node[box, fill=gray!15] (in_write) {
    \textbf{From Core $i$ (100 MHz)}\\
    \scriptsize\texttt{x\_w[8:0]}, \texttt{y\_w[8:0]}\\
    \scriptsize\texttt{data\_w[14:0]} (RGB555)\\
    \scriptsize\texttt{z\_w[3:0]} (Depth)\\
    \scriptsize\texttt{enb} (Pixel Valid)
};

% ==================== ADDRESS DECOMPOSITION ====================
\node[logicbox, right=1.0cm of in_write, yshift=1.2cm] (addr_math) {
    \textbf{Multiplier-Free Address}\\
    \scriptsize $Y_{\text{local}} = Y_w - (i \times 24)$\\
    \scriptsize $\text{Addr} = (Y_{\text{local}} \ll 8) + (Y_{\text{local}} \ll 6) + X_w$\\
    \scriptsize No DSP48 Blocks Required!
};
\draw[bus] (in_write.north) |- (addr_math.west);

% ==================== Z-BUFFER UNIT ====================
\node[logicbox, fill=red!10, draw=red!70, right=1.0cm of in_write, yshift=-1.2cm] (z_unit) {
    \textbf{Hardware Z-Buffer Comparator}\\
    \scriptsize Read $Z_{\text{mem}} = \text{ZRAM}[\text{Addr}]$\\
    \scriptsize \texttt{z\_pass <= (z\_w <= z\_mem)}\\
    \scriptsize Updates ZRAM on success
};
\draw[bus] (in_write.south) |- (z_unit.west);

% ==================== VRAM & ZRAM DUAL-PORT ====================
\node[membox, right=1.4cm of addr_math] (vram) {
    \textbf{Color VRAM Dual-Port}\\
    \scriptsize $2 \times 7680 = 15360$ words $\times$ 15-bit\\
    \scriptsize \textbf{Port A (Write)}: 100 MHz\\
    \scriptsize \textbf{Port B (Read)}: 25 MHz
};

\node[membox, fill=red!5, draw=red!60, right=1.4cm of z_unit] (zram) {
    \textbf{Depth ZRAM}\\
    \scriptsize $7680 \times 4$-bit\\
    \scriptsize Reset to \texttt{"1111"} (Furthest)
};

\draw[bus] (addr_math.east) -- node[above, font=\tiny] {Write Addr} (vram.west);
\draw[bus] (addr_math.east) -- (zram.west);
\draw[ctl] (z_unit.east) -- node[below, font=\tiny] {\texttt{write\_enb = enb AND z\_pass}} (vram.210);

% ==================== FSM & SWAP LOGIC ====================
\node[fsmbox, below=1.0cm of z_unit, xshift=1.5cm] (swap_fsm) {
    \textbf{Swap \& Auto-Clear FSM}\\
    \scriptsize State 0: Normal Drawing (Back Buffer)\\
    \scriptsize State 1: Wait for \texttt{v\_sync = 0} (Blanking)\\
    \scriptsize State 2: Hardware Auto-Clear to White (0x7FFF) \& $Z=15$\\
    \scriptsize State 3: Swap Complete (\texttt{swapped = 1})
};

\draw[ctl] (swap_fsm.north) -- node[right, font=\tiny] {Buffer Inversion} (vram.south);
\draw[ctl] (swap_fsm.north) -- (zram.south);

% ==================== DISPLAY SCANNING PORT B ====================
\node[box, fill=blue!10, draw=blue!60, right=1.5cm of vram] (disp_port) {
    \textbf{To Display Controller}\\
    \scriptsize\texttt{clock\_r} = 25 MHz\\
    \scriptsize\texttt{x\_r[8:0]}, \texttt{y\_r[8:0]}\\
    \scriptsize\texttt{data\_r[14:0]} (Front Buffer)
};
\draw[bus] (vram.east) -- (disp_port.west);

\end{tikzpicture}

}
\caption{Microarchitecture of the Tile Framebuffer and Integrated Z-Buffer (\texttt{tile.vhd}).}
\label{fig:tile_memory}
\end{figure}

\section{True Dual-Port BRAM Organization and Double Buffering}
Each tile canvas has a resolution of $320 \times 24 = 7{,}680$ pixels. To support double buffering, the color VRAM stores two complete frames:
\begin{equation}
\text{VRAM Size} = 2 \times (320 \times 24) = 15{,}360 \text{ words} \times 15 \text{ bits}
\end{equation}

The ping-pong register \texttt{current\_fb} controls the base address offsets for read and write ports:
\begin{itemize}
    \item When \texttt{current\_fb = '0'}: The compute core writes into the \textbf{Back Buffer} (starting at address offset $7{,}680$), while the video display scans from the \textbf{Front Buffer} (starting at address offset $0$).
    \item When \texttt{current\_fb = '1'}: The compute core writes into the \textbf{Back Buffer} (offset $0$), while the video display scans from the \textbf{Front Buffer} (offset $7{,}680$).
\end{itemize}

\section{Swap Protocol and Hardware Background Auto-Clear}
The buffer swapping sequence is coordinated through a 4-state FSM in each tile:
\begin{enumerate}
    \item \textbf{State 0 (Drawing)}: Normal rendering into the back buffer. The tile monitors \texttt{fb\_swap}.
    \item \textbf{State 1 (Wait for VSYNC)}: When all 10 cores have processed the \texttt{SWAP\_BUFFERS} instruction, the global signal \texttt{fb\_swap\_general} goes high. The FSM waits for the vertical blanking interval of the display (\texttt{v\_sync\_sync = '0'}). Once detected, \texttt{current\_fb} is toggled:
    \begin{equation}
    \texttt{current\_fb} \gets \neg \texttt{current\_fb}
    \end{equation}
    Because the pointer inversion occurs strictly during vertical blanking, screen tearing is completely eliminated.
    \item \textbf{State 2 (Hardware Auto-Clear)}: During this state, an internal counter \texttt{clear\_cnt} increments from $0$ to $7{,}679$, sweeping the newly assigned back buffer. At each cycle, the hardware writes white ($15\text{-bit } \texttt{"111111111111111"}$) into VRAM and maximum distance ($\texttt{"1111"}$) into ZRAM.
    \item \textbf{State 3 (Swap Complete)}: Asserts \texttt{swapped <= '1'} to inform \texttt{spEXEC} that the new buffer is clean and ready for rendering.
\end{enumerate}

\section{Integrated 4-Bit Hardware Z-Buffer (Depth Test)}
To enable proper 2.5D and 3D visual layering without requiring the host CPU to sort polygons from back-to-front (Painter's Algorithm), each tile memory module embeds a dedicated 4-bit depth memory (\texttt{ZRAM}) of $7{,}680$ entries.

\subsection{Single-Cycle Depth Comparator}
On every clock cycle where a candidate pixel is asserted (\texttt{enb = '1'}), the hardware performs an instantaneous read-compare-write cycle:
\begin{equation}
Z_{\text{current}} = \text{ZRAM}[\text{address\_w}]
\end{equation}
\begin{equation}
\texttt{z\_pass} = \begin{cases} 
\text{'1'} & \text{if } Z_{\text{in}} \le Z_{\text{current}} \\ 
\text{'0'} & \text{otherwise} 
\end{cases}
\end{equation}
\begin{equation}
\texttt{write\_enb} = \texttt{enb} \land \texttt{hit\_w} \land \texttt{z\_pass}
\end{equation}

If the candidate pixel's depth $Z_{\text{in}}$ is less than or equal to the current depth stored in \texttt{ZRAM} (indicating that the new primitive is closer to or coplanar with the viewer), \texttt{write\_enb} is asserted: the pixel color is committed to VRAM, and \texttt{ZRAM} is updated with $Z_{\text{in}}$. Otherwise, the write is suppressed, effectively culling obscured fragments at zero performance penalty.

\section{Multiplier-Free Address Generation}
In standard architectures, mapping 2D coordinates $(X, Y)$ to a linear 1D memory array requires integer multiplication:
\begin{equation}
\text{Addr} = Y_{\text{local}} \times 320 + X
\end{equation}
Implementing ten such multipliers in parallel would consume 10 to 20 DSP48 slices solely for address calculation. 

spGPU circumvents this by decomposing $320$ into powers of two ($320 = 256 + 64 = 2^8 + 2^6$):
\begin{equation}
\text{Addr} = (Y_{\text{local}} \ll 8) + (Y_{\text{local}} \ll 6) + X
\end{equation}

Similarly, the tile vertical offset $Y_{\text{offset}} = \text{tile\_index} \times 24$ is decomposed as $24 = 16 + 8 = 2^4 + 2^3$:
\begin{equation}
Y_{\text{offset}} = (\text{tile\_index} \ll 4) + (\text{tile\_index} \ll 3)
\end{equation}
As shown in Listing~\ref{lst:tile_addr_math}, this enables the synthesizer to map address generation entirely to barrel shifters and carry-chain adders, conserving all FPGA DSP48 blocks for geometry and filtering.

\begin{lstlisting}[language=VHDL, caption={Shift-and-Add Memory Address Generation in \texttt{tile.vhd}.}, label={lst:tile_addr_math}]
-- Offset for the tile: tile_index * 24 = (tile_index * 16) + (tile_index * 8)
y_offset <= shift_left(resize(unsigned(tile_index), N_PIXEL), 4) + 
            shift_left(resize(unsigned(tile_index), N_PIXEL), 3);

-- Local address calculation: Y_local * 320 + X = (Y_local * 256) + (Y_local * 64) + X
address_w <= std_logic_vector(
    shift_left(resize(y_w_local, 18), 8) + 
    shift_left(resize(y_w_local, 18), 6) + 
    resize(unsigned(x_w), 18)
);
\end{lstlisting}
